\begin{definition}
	Пусть $E$ "--- линейное нормированное пространство. Система $\{e_n\} \subset E$ называется \textit{базисом (Шаудера)} в $E$, если для любого $x \in E$ существует единственный набор чисел $\{\alpha_n\} \subset \K$ такой, что 
    $x = \sum \limits_{n = 1}^{\infty} \alpha_n e_n$.
\end{definition}

\begin{proposition}[неравенство Бесселя]
	Пусть $E$ "--- евклидово пространство, элементы 
    $\{e_i\}$ из $E$ образуют ортонормированную
    систему. Тогда для любого $x \in E$ выполнено следующее
    неравенство:
	\[\sum_{k = 1}^{\infty}|(x, e_k)|^2 \leqslant \|x\|^2\]
\end{proposition}

\begin{proof}
	Достаточно заметить, что в силу ортнормированности системы $\{e_k\}$ выполнено следующее равенство:
	\[ 0 \leqslant \left\|x - \sum_{k = 1}^n(x, e_k)e_k\right\|^2 = 
 \|x\|^2 - \sum_{k = 1}^{n}|(x, e_k)|^2\]

	Поскольку левая часть равенства выше неотрицательна, то неотрицательна и правая часть, и верно это для произвольного $n \geqslant 1$ из чего следует требуемое.
\end{proof}

\begin{corollary}[неравенство Коши---Буняковского]
	Пусть $E$ "--- евклидово пространство. Тогда для любых $x, y \in E$ выполнено следующее неравенство:
	\[|(x, y)| \leqslant \|x\| \cdot \|y\|\]
\end{corollary}

\begin{proof}
	Случай, когда $y = 0$, тривиален, поэтому можно считать, что $y \neq 0$. Тогда система $\left\{\frac{y}{\|y\|}\right\}$ является ортонормированной, и по неравенству Бесселя выполнено следующее:
	\[\left|\left(x, \frac y{\|y\|}\right)\right|^2 \leqslant \|x\|^2 \ra |(x, y)| \leqslant \|x\| \cdot \|y\| \qedhere\]
\end{proof}

\begin{note}
	Из неравенства Коши-Буняковского, в частности, следует, что скалярное произведение непрерывно на $E$ относительно порождаемой им нормы.
\end{note}

\begin{proposition}\label{minproperty}
	Пусть $E$ "--- евклидово пространство, элементы $e_1, \dotsc, e_k \in E$ образуют ортонормированную систему. Тогда для любого $x \in E$ и любого набора $\{\alpha_k\}_{k = 1}^n \subset \K$ выполнено следующее неравенство:
	\[\left\|x - \sum_{k = 1}^n \alpha_ke_k\right\| \geqslant \left\|x - \sum_{k = 1}^n (x, e_k)e_k\right\|\]

	Более того, равенство в неравенстве выше достигается тогда и только тогда, когда $\alpha_k = (x, e_k)$ для всех $k \in \{1, \dotsc, n\}$.
\end{proposition}

\begin{proof}[Доказательство для случая, когда {$\K = \R$}]
	В силу ортонормированности системы $\{e_k\}_{k = 1}^n$, выполнены следующие равенства:
	\begin{multline*}
		\left\|x - \sum_{k = 1}^n \alpha_ke_k\right\|^2
		=
		\left(x - \sum_{k = 1}^n \alpha_ke_k, x - \sum_{k = 1}^n \alpha_ke_k\right)
		= 
		\\
		=
		\|x\|^2 - 2\sum_{k = 1}^n\alpha_k(x, e_k) + \sum_{k = 1}^n\alpha_k^2 = \left\|x - \sum_{k = 1}^n(x, e_k)e_k\right\|^2 + \sum_{k = 1}^n\big((x, e_k) - \alpha_k\big)^2
	\end{multline*}

	Из равенства между самой левой и самой правой частями в цепочке выше получаем требуемое.
\end{proof}

\begin{theorem}\label{thm4.4}
	Пусть $H$ "--- сепарабельное гильбертово пространство, $\dim{H} = +\infty$, и ${e = \{e_n\} }$ "--- ортонормированная система. Тогда следующие условия эквивалентны:
	\begin{enumerate}
		\item $e$ "--- ортонормированный базис в $H$
  		\item $\overline{[e]} = H$ т.е. $e$ -- полная система
    	\item Для любого $h \in H$ выполнено равенство Парсеваля: $\|h\|^2 = \sum \limits_{n = 1}^\infty |(h, e_n)|^2$
    	\item $e^\perp = \{0\}$
	\end{enumerate}
\end{theorem}

\begin{proof}~
	\begin{itemize}
		\item\imp{1}{2}Очевидно из определения базиса.
  
		\item\imp{2}{1}Зафиксируем произвольный $h \in H$ и произвольное $\varepsilon > 0$. По условию, существует конечный набор $\alpha_1, \dotsc, \alpha_n$ такой, что $\|h - \sum_{k = 1}^n\alpha_k e_k\| < \varepsilon$. Тогда, по утверждению \ref{minproperty}, выполнено также неравенство $\|h - \sum_{k = 1}^n(h, e_k) e_k\| < \varepsilon$. Тогда, в силу произвольности числа $\varepsilon$, выполнено равенство $h = \sum_{n = 1}^\infty (h, e_n)e_n$.
		
		Проверим теперь, что разложение элемента $h$ единственно. Пусть для некоторого набора $\{\beta_n\} \subset \K$ выполнено равенство $h = \sum_{n = 1}^\infty \beta_n e_n$. Тогда для любого $k \in \N$, скалярно умножая частичную сумму ряда $\sum_{n = 1}^\infty\beta_n e_n$ на $e_k$ и переходя к пределу, получаем, что $\beta_k = (h, e_k)$, что и означает требуемое.
  
		\item\eqv{1}{3}Уже было замечено, что для любого $h \in H$ и любого $k \in \N$ выполнено следующее равенство:
		\[\left\|h - \sum_{k = 1}^n(h, e_k)e_k\right\|^2 = \|h\|^2 - \sum_{k = 1}^n|(h, e_k)|^2\]
		
		Значит, $h = \sum_{n = 1}^\infty (h, e_n)e_n \lra \|h\|^2 = \sum_{n = 1}^\infty |(h, e_n)|^2$, и единственность разложения элемента $h$ доказывается так же, как в импликации выше.
  
		\item\eqv{2}{4}Из теоремы Рисса  и равенства $e^{\perp} = \overline{[e]}^{\perp} = M^{\perp}$ получаем требуемое. То есть \\ если система полна, то $\overline{[e]} = H \; \Rightarrow \; M^{\perp} = \{0\}$. Если $e^\perp = \{0\} \; \Rightarrow \; M = H $, то система полна. \qedhere
	\end{itemize}
\end{proof}

\begin{corollary}
	Пусть $H$ "--- сепарабельное гильбертово пространство. Тогда в $H$ существует не более чем счётный ортонормированный базис.
\end{corollary}

\begin{proof}
	Случай, когда $\dim H <+\infty$, уже рассматривался в курсе алгебры, поэтому будем считать, что $\dim H = +\infty$. Рассмотрим счётное множество $F \subset H$ такое, что $\overline F = H$(оно существует по определению сепарабельности \ref{separability}), тогда $F = \{f_n\}$. Проведём процесс ортогонализации Грама--Шмидта и получим ортонормированную систему $e$, для которой выполнено $[e] = [F]$, откуда $\overline{[e]} = H$, тогда, по теореме \ref{thm4.4}, $e$ "--- ортонормированный базис в $H$.  
\end{proof}

\begin{note}
	Результат выше позволяет построить изоморфизм между любыми двумя сепарабельными гильбертовыми пространствами бесконечной размерности: достаточно отобразить ортонормированный базис в одном пространстве в ортонормированный базис в другом пространстве.
\end{note}
